% Generated draft: requires mathematical and source review.
\documentclass[11pt]{article}
\usepackage[margin=25mm]{geometry}
\usepackage{fontspec}
\setmainfont{Latin Modern Roman}
\usepackage{amsmath,amssymb,mathtools,unicode-math}
\setmathfont{Latin Modern Math}
\usepackage{enumitem,needspace,xurl,hyperref}
\setlength{\parindent}{0pt}
\setlength{\parskip}{5pt}
\setlength{\emergencystretch}{3em}
\newcommand{\N}{\mathbb{N}}
\newcommand{\Z}{\mathbb{Z}}
\newcommand{\Q}{\mathbb{Q}}
\newcommand{\R}{\mathbb{R}}
\newcommand{\C}{\mathbb{C}}
\newcommand{\F}{\mathbb{F}}
\newcommand{\A}{\mathbb{A}}
\newcommand{\PP}{\mathbb{P}}
\newcommand{\E}{\mathbb{E}}
\newcommand{\eps}{\varepsilon}
\newcommand{\dd}{\,\mathrm{d}}
\begin{document}
\section*{UnsolvedMath ID: 30002782}
\textbf{Source problem code:} catalogue entry
\subsection{Definitions and mathematical statement}
For a graph \(G=(V,E)\), its independence complex \(\operatorname{Ind}(G)\) is the abstract simplicial complex whose simplices are the finite independent subsets of \(V\), meaning subsets containing no edge of \(G\). The source does not specify finiteness or graph conventions, nor the coefficient ring or whether cohomology is reduced.

\textbf{Edge-replacement problem.} Let \(G\) be a graph for which \(\operatorname{Ind}(G)\) has, in the source's words, ``exactly one cohomology generator,'' and let \(e\) be an edge of \(G\). For an edge with distinct endpoints \(u,v\), replacing \(e\) by a path of \(3k+1\) edges means deleting \(e\), introducing \(3k\) new internal vertices, and joining \(u\) to \(v\) through those vertices, while retaining the other vertices and edges. Denote the resulting graph by \(G_k\). This construction makes sense for integers \(k\geq 0\); the source does not specify whether the requested family includes \(k=0\).

For each admissible pair \((G,e)\), find a closed formula and a natural generating function, as \(k\) varies, for the quantity described in the source as ``the largest weight of the cohomology representatives'' of \(\operatorname{Ind}(G_k)\). The supplied record does not define weight, the allowed representatives, their normalization, or the maximization being requested. Moreover, ``exactly one cohomology generator'' does not specify a cohomology theory, degree, or coefficient structure. Consequently, this is a qualified statement of the requested target, not a definition of a numerical invariant; those missing data must be supplied before a precise formula or generating function can be sought.

\textbf{Alternating-sign-matrix realization problem.} An alternating-sign matrix of order \(n\geq 1\) is an \(n\times n\) matrix with entries in \(\{-1,0,1\}\), with each row and column summing to \(1\), and with the nonzero entries in every row and column alternating in sign. Find independence complexes realizing the symmetry classes or versions of alternating-sign matrices intended by the source's reference to Kuperberg. The supplied record neither enumerates these classes nor defines ``realizing'': it does not specify a bijection, an enumerative identity, or a cohomological construction. These choices remain unresolved and must not be inferred from the word ``realizing.''
\subsection{Short English statement}
The first question asks how a specified cohomological weight changes when an edge is replaced by a path of \(3k+1\) edges. The second asks for independence complexes corresponding to Kuperberg's alternating-sign-matrix symmetry classes. The supplied record leaves the weight, cohomology conventions, and intended correspondence undefined, so both targets require clarification before a fully precise formulation is possible.
\subsection{Sources}
\begin{itemize}
\item Supplied local UnsolvedMath registry record, permanent ID \(30002782\), title \emph{Independence Complexes and Alternating-Sign Matrices}, exact source problem code \texttt{catalogue entry}, in the Oberwolfach Reports Open Problems collection. The record gives proposed year \(2015\) and supplies no named proposer.
\item \emph{Geometric and Algebraic Combinatorics} (2015), Oberwolfach Reports, pp.\ 285--368, DOI \texttt{10.4171/owr/2015/5}. Supplied source report and literature-source link: \url{https://doi.org/10.4171/owr/2015/5}. The supplied background describes Oberwolfach Reports as documenting mathematical workshops and their open-problem discussions.
\item Kuperberg, symmetry classes of alternating-sign matrices, supplied literature-source link: \url{https://doi.org/10.2307/3597283}. The record supplies no more detailed bibliographic description.
\item Supplied literature assessment dated \texttt{2026-08-21}: the record labels the problem \texttt{partially\_solved} and states that cycle and path families have explicit alternating-sign-matrix sequences and suspension recurrences, but that no general closed formula for every graph-edge replacement or realization of all Kuperberg symmetry classes was verified. This is supplied dated curation derived from the OWR material, not an independently verified literature assessment. The background requests independent verification of time-sensitive status and cited claims.
\item The registry labels the statement \texttt{corrected\_verified} and reports that the official Oberwolfach passage was checked directly and its mathematical question, definitions, and notation restored without adding a stronger claim. This verification assertion is supplied provenance; no independent source consultation is asserted here.
\end{itemize}
\end{document}
